\subsection{PASSAGE FROM LIE ALGEBRAS TO LIE GROUPS}

THEOREM 3. (i) If L is a finite-dimensional Lie algebra, there exists a simply connected Lie group G such that L(G) is isomorphic to L.

\begin{enumerate}
\def\labelenumi{(\roman{enumi})}
\setcounter{enumi}{1}
\tightlist
\item
  Let \(G_{1}\) and \(G_{2}\) be two connected Lie groups, with \(G_{1}\) simply connected. Let f be an isomorphism of \(\mathrm{L}(\mathrm{G}_{1})\) onto \(\mathrm{L}(\mathrm{G}_{2})\) , \(\phi\) the morphism of \(G_{1}\) into \(G_{2}\) such that \(\mathrm{L}(\phi)=f\) and N the kernel of \(\phi\) . Then N is a discrete subgroup of the centre of \(G_{1}\) and the morphism of \(G_{1}/N\) into \(G_{2}\) derived from \(\phi\) is a Lie group isomorphism. If \(G_{2}\) is simply connected, \(\phi\) is an isomorphism.
\end{enumerate}

Let \(\mathbf{L}\) be a finite-dimensional Lie algebra. There exists a finite-dimensional vector space \(\mathbf{E}\) such that \(\mathbf{L}\) can be identified with a Lie subalgebra of \(\operatorname{End}(\mathbf{E})\) (Chapter I, §7, Theorem 2). Let \(\mathbf{H}\) be the integral subgroup of \(\mathbf{GL}(\mathbf{E})\) with

Lie algebra L. Let \(\hat{\mathbf{H}}\) be its universal covering (§ 1, no. 9, Remark). Then \(\mathbf{L}(\hat{\mathbf{H}})\) is isomorphic to L, whence (i).

Let \(G_{1}, G_{2}, \tilde{f}, \phi, N\) be as in (ii). Then \(\phi\) is étale, hence \(\phi(G_{1})\) is an open subgroup of \(G_{2}\) and hence \(\phi(G_{1}) = G_{2}\) . On the other hand, N is discrete and hence contained in the centre of \(G_{1}\) (Integration, Chapter VII, § 3, Lemma 4). Clearly the morphism of \(G_{1}/N\) onto \(G_{2}\) derived from \(\phi\) is a Lie group isomorphism. If \(G_{2}\) is simply connected, every étale mapping of \(G_{1}\) onto \(G_{2}\) is injective and hence \(N = \{e\}\) .

PROPOSITION 5. Let G be a connected real Lie group. Suppose that L(G) has a complex normable Lie algebra structure L' compatible with its real normable Lie algebra structure. There exists on G one and only one complex Lie group structure compatible with the real Lie group structure and with Lie algebra L'.

By § 4, no. 2, Corollary 2 to Theorem 2, it suffices to prove that the structure of \(L'\) is invariant under Ad G. Let \(\phi\) be an exponential mapping of G. By § 4, no. 4, Corollary 3 (i) to Proposition 8, there exists a neighbourhood V of 0 in \(L(G)\) such that the structure of \(L'\) is invariant under \(Ad\phi(V)\) . But \(Ad\phi(V)\) generates G because G is connected.

The conclusion of Proposition 5 is not necessarily true if G is not assumed to be connected (Exercise 7).

PROPOSITION 6. Let G be a connected complex Lie group. If G is compact, G is commutative.

The holomorphic mapping \(g \mapsto \operatorname{Ad} g\) of \(\mathbf{G}\) into \(\mathcal{L}(\mathbf{L}(\mathbf{G}))\) is constant (Differentiable and Analytic Manifolds, R, 3.3.7) and hence ad \(a = 0\) for all \(a \in \mathbf{L}(\mathbf{G})\) (§ 3, no. 12, Proposition 44). Hence \(\mathbf{G}\) is commutative (§ 4, Corollary 3 to Theorem 1).

